\documentclass[11pt, letterpaper]{article}
\input{../../homework.tex}
\usepackage{titlepageBU}
\title{Homework 1}
\courseID{MA 542}
\courseSection{A1}
\begin{document}
\makeproblem
\section*{Problem 1}
The set $\left\{ (1,0),(0,1) \right\} $ forms a basis for $\mathbb{R} ^2$ as an $\mathbb{R} $-vector space. To prove this, let $(x,y)\in\mathbb{R}\times \mathbb{R}  $. Since $x,y\in\mathbb{R} $,
\[
	(x,y)=x(1,0)+y(0,1)
.\]
Also note that the two elements are trivially linearly independent, so they form a basis. Consider the map $\phi :\mathbb{R} \times \mathbb{R} \to \mathbb{C} $ defined by
\[
	(a,b)\mapsto a+bi
.\]
Let $r\in\mathbb{R} $ and $(c,d)\in\mathbb{R} \times \mathbb{R} $, and note that
\[
	\phi (r(a,b))=r(a+bi)=r\phi (a,b),
\]
\[
	\phi ((a,b)+(c,d))=\phi (a+c,b+d)=(a+c)+(b+d)i=(a+bi)+(c+di)=\phi (a,b)+\phi (c,d)
.\]
Therefore, $\phi $ is an $\mathbb{R} $-linear map, and it also gives a bijection of basis elements:
\[
	\phi (1,0)=1,\qquad \phi (0,1)=i
.\]
Therefore, $\phi $ is a vector space isomorphism, and $\mathbb{R} \times \mathbb{R} \cong \mathbb{C} $.
\section*{Problem 2}
No to both. Note that $(1,0),(0,1)\neq 0\coloneqq (0,0)$. Also note that
\[
	(1,0)\cdot (0,1)=(1\cdot 0,0\cdot 1)=(0,0)
.\]
Therefore, $\mathbb{R} \times \mathbb{R} $ has zero divisors, and is not an integral domain nor a field.
\section*{Problem 3}
No. The number $i$ must satisfy $i^2=-1$. In $\mathbb{R} \times \mathbb{R} $, $-1$ is simply the additive inverse of $1$, and $1=(1,1)$ fairly obviously:
\[
	(a,b)(1,1)=(a1,b1)=(a,b)
.\]
We thus seek an element $(a,b)\in\mathbb{R}^2 $ such that $(a,b)^2=(-1,-1)$. Since multiplication is defined elementwise, it follows that $a^2,b^2=-1$, with $a,b\in\mathbb{R} $. This is a contradiction because there exists no $a\in\mathbb{R} $ such that $a^2=-1$.
\end{document}
